% GENERATED FILE. Edit canonical lessons, then run npm run generate:books.
% canonical-source-hash: fnv1a64:436b7c39035865a3
% canonical-lessons: AR-C260-rashsha, AR-C260-taraqa, AR-C260-alsaqa, AR-C260-harraka, AR-C260-laffa

\chapter{To sprinkle, To knock, To stick, To stir, To wrap}
\label{ch:ar-to-sprinkle-to-knock-to-stick-to-stir-to-wrap}
\hlchaptermodality{260}

\begin{chapteropening}
\textbf{By the end of this chapter:} \emph{I can say to sprinkle, to knock, to stick, to stir and to wrap.}

Complete the last lesson of chapter 260: I can say to sprinkle, to knock, to stick, to stir and to wrap.
\end{chapteropening}

\section[\texorpdfstring{\ar{رش} (rashsha)}{رش (rashsha)}]{\ar{رش} (rashsha) — to sprinkle, to spray}

\label{lesson:AR-C260-rashsha}



\begin{quote}
Before the new one: say the Arabic for to jump, then the Arabic for to press.
\end{quote}

\subsection*{You'll want to know: \ar{رش}}
\textbf{\ar{رش}} — \emph{rashsha} — \textquotedblleft{}to sprinkle, to spray\textquotedblright{}.

\begin{cousinweb}
\textbf{\ar{رش}} (\emph{rashsha}) is built on the native Arabic root \emph{r-sh-sh}, \textquotedblleft{}to sprinkle\textquotedblright{} (Lane). \emph{Rashsha} is the plain verb, \textquotedblleft{}he sprinkled\textquotedblright{}.
\end{cousinweb}

\begin{grammarlens}[title={What you've built}]
One more word for this chapter.
\end{grammarlens}

\subsection*{Your turn}
\begin{itemize}
\raggedright
  \item \emph{Say it:} \emph{rashsha}
  \item \emph{Say it:} \emph{rashsha}, once more
  \item \emph{Say it:} say \emph{rashsha}
  \item \emph{Recall:} say the Arabic for to jump, then the Arabic for to press, then say \emph{rashsha} again
\end{itemize}

\begin{tcolorbox}[breakable,colback=teal!4,colframe=teal!35!black,title={Before you move on}]
What does \ar{رش} mean? (\textquotedblleft{}To sprinkle, to spray\textquotedblright{}.) Say it once more.
\end{tcolorbox}
\section[\texorpdfstring{\ar{طرق} (ṭaraqa)}{طرق (ṭaraqa)}]{\ar{طرق} (ṭaraqa) — to knock (on a door)}

\label{lesson:AR-C260-taraqa}



\begin{quote}
Before the new one: say the Arabic for to press, then the Arabic for to sprinkle.
\end{quote}

\subsection*{You'll want to know: \ar{طرق}}
\textbf{\ar{طرق}} — \emph{ṭaraqa} — \textquotedblleft{}to knock (on a door)\textquotedblright{}.

\begin{cousinweb}
\textbf{\ar{طرق}} (\emph{ṭaraqa}) is built on the native Arabic root \emph{ṭ-r-q}, \textquotedblleft{}to strike, to knock\textquotedblright{} (Lane). \emph{Ṭaraqa} is the plain verb, \textquotedblleft{}he knocked\textquotedblright{}.
\end{cousinweb}

\begin{grammarlens}[title={What you've built}]
One more word for this chapter.
\end{grammarlens}

\subsection*{Your turn}
\begin{itemize}
\raggedright
  \item \emph{Say it:} \emph{ṭaraqa}
  \item \emph{Say it:} \emph{ṭaraqa}, once more
  \item \emph{Say it:} say \emph{ṭaraqa}
  \item \emph{Recall:} say the Arabic for to press, then the Arabic for to sprinkle, then say \emph{ṭaraqa} again
\end{itemize}

\begin{tcolorbox}[breakable,colback=teal!4,colframe=teal!35!black,title={Before you move on}]
What does \ar{طرق} mean? (\textquotedblleft{}To knock (on a door)\textquotedblright{}.) Say it once more.
\end{tcolorbox}
\section[\texorpdfstring{\ar{ألصق} (alṣaqa)}{ألصق (alṣaqa)}]{\ar{ألصق} (alṣaqa) — to stick, to glue}

\label{lesson:AR-C260-alsaqa}



\begin{quote}
Before the new one: say the Arabic for to sprinkle, then the Arabic for to knock.
\end{quote}

\subsection*{You'll want to know: \ar{ألصق}}
\textbf{\ar{ألصق}} — \emph{alṣaqa} — \textquotedblleft{}to stick, to glue\textquotedblright{}.

\begin{cousinweb}
\textbf{\ar{ألصق}} (\emph{alṣaqa}) is built on the native Arabic root \emph{l-ṣ-q}, \textquotedblleft{}to cleave, to stick\textquotedblright{} (Lane). \emph{Alṣaqa} is Form IV, \textquotedblleft{}to make stick\textquotedblright{}.
\end{cousinweb}

\begin{grammarlens}[title={What you've built}]
One more word for this chapter.
\end{grammarlens}

\subsection*{Your turn}
\begin{itemize}
\raggedright
  \item \emph{Say it:} \emph{alṣaqa}
  \item \emph{Say it:} \emph{alṣaqa}, once more
  \item \emph{Say it:} say \emph{alṣaqa}
  \item \emph{Recall:} say the Arabic for to sprinkle, then the Arabic for to knock, then say \emph{alṣaqa} again
\end{itemize}

\begin{tcolorbox}[breakable,colback=teal!4,colframe=teal!35!black,title={Before you move on}]
What does \ar{ألصق} mean? (\textquotedblleft{}To stick, to glue\textquotedblright{}.) Say it once more.
\end{tcolorbox}
\section[\texorpdfstring{\ar{حرك} (ḥarraka)}{حرك (ḥarraka)}]{\ar{حرك} (ḥarraka) — to stir, to move (something)}

\label{lesson:AR-C260-harraka}



\begin{quote}
Before the new one: say the Arabic for to knock, then the Arabic for to stick.
\end{quote}

\subsection*{You'll want to know: \ar{حرك}}
\textbf{\ar{حرك}} — \emph{ḥarraka} — \textquotedblleft{}to stir, to move (something)\textquotedblright{}.

\begin{cousinweb}
\textbf{\ar{حرك}} (\emph{ḥarraka}) is built on the native Arabic root \emph{ḥ-r-k}, \textquotedblleft{}to move\textquotedblright{} (Lane). \emph{Ḥarraka} is Form II, \textquotedblleft{}to set moving\textquotedblright{}.
\end{cousinweb}

\begin{grammarlens}[title={What you've built}]
One more word for this chapter.
\end{grammarlens}

\subsection*{Your turn}
\begin{itemize}
\raggedright
  \item \emph{Say it:} \emph{ḥarraka}
  \item \emph{Say it:} \emph{ḥarraka}, once more
  \item \emph{Say it:} say \emph{ḥarraka}
  \item \emph{Recall:} say the Arabic for to knock, then the Arabic for to stick, then say \emph{ḥarraka} again
\end{itemize}

\begin{tcolorbox}[breakable,colback=teal!4,colframe=teal!35!black,title={Before you move on}]
What does \ar{حرك} mean? (\textquotedblleft{}To stir, to move (something)\textquotedblright{}.) Say it once more.
\end{tcolorbox}
\section[\texorpdfstring{\ar{لف} (laffa)}{لف (laffa)}]{\ar{لف} (laffa) — to wrap, to roll up}

\label{lesson:AR-C260-laffa}



\begin{quote}
Before the new one: say the Arabic for to stick, then the Arabic for to stir.
\end{quote}

\subsection*{You'll want to know: \ar{لف}}
\textbf{\ar{لف}} — \emph{laffa} — \textquotedblleft{}to wrap, to roll up\textquotedblright{}.

\begin{cousinweb}
\textbf{\ar{لف}} (\emph{laffa}) is built on the native Arabic root \emph{l-f-f}, \textquotedblleft{}to wrap, to roll up\textquotedblright{} (Lane). \emph{Laffa} is the plain verb, \textquotedblleft{}he wrapped\textquotedblright{}.
\end{cousinweb}

\begin{grammarlens}[title={What you've built}]
One more word for this chapter.
\end{grammarlens}

\subsection*{Your turn}
\begin{itemize}
\raggedright
  \item \emph{Say it:} \emph{laffa}
  \item \emph{Say it:} \emph{laffa}, once more
  \item \emph{Say it:} say \emph{laffa}
  \item \emph{Recall:} say the Arabic for to stick, then the Arabic for to stir, then say \emph{laffa} again
\end{itemize}

\begin{tcolorbox}[breakable,colback=teal!4,colframe=teal!35!black,title={Before you move on}]
What does \ar{لف} mean? (\textquotedblleft{}To wrap, to roll up\textquotedblright{}.) Say it once more.
\end{tcolorbox}
